\begin{textsample}{POS Dim 1 – 2012\_c5 – Score 6.00 – 2012\_c5/enade\_2012\_jorn\_d0855fab9c.txt}  \label{ex:f1_pos_115}
CAPÍTULO 5 DISTRIBUIÇÃO DOS CONCEITOS 5.1 \textbf{PANORAMA} NACIONAL DA DISTRIBUIÇÃO DOS CONCEITOS A Tabela 5.1 apresenta a quantidade e distribuição de cursos de Comunicação Social - Jornalismo participantes do ENADE/2012, por faixa de conceito e Grande Região.
A diferença entre os cursos \textbf{tabulados} neste capítulo e no capítulo 2 são os cursos sem conceito, em princípio, aqueles sem alunos concluintes que participassem da prova ou que tivessem tão somente um aluno concluinte¹¹.
Observando -se os dados da Tabela 5.1, nota -se que, dos 277 cursos participantes, 119 ( 43,0 \% ) classificaram -se com conceito 3, o valor modal.
Este foi também o conceito modal em todas as regiões do Brasil ( entre 39,3 \% no Sudeste e 55,1 \% no Nordeste ), exceto a região Sul, onde o conceito modal foi 4 ( 38,2 \% ).
O conceito 4 foi o segundo mais frequente em nível nacional ( 35,0 \%, correspondendo a 97 cursos ) e o conceito 2, o terceiro ( 11,6 \%, correspondendo a 32 cursos ).
Houve, ainda, 15 cursos ( 5,4 \% ) que receberam conceito 5 e outros nove cursos que receberam conceito 1 ( 3,2 \% ).
Dos 277 cursos de Comunicação Social – Jornalismo, cinco ficaram sem conceito ( SC ).
A região Norte participou com 23 cursos ou 8,3 \% do total nacional.
Nesta região, o conceito 3 foi atribuído a dez cursos dos 23 participantes, o que \textbf{equivale} a 43,5 \% do total regional.
Nenhum curso ficou sem conceito ( SC ).
Os demais foram avaliados com conceito 1 ( um curso, correspondendo a 4,3 \% ), conceito 2 ( quatro cursos correspondendo a 17,4 \% ) e conceito 4 ( oito cursos correspondendo a 34,8 \% ).
Nenhum curso foi alocado ao conceito 5. ¹¹ Estes cursos com somente um concluinte foram considerados no capítulo 2.
A região Nordeste participou com 49 cursos ou 17,7 \% do total nacional.
Como já comentado, destes, 27 cursos, 55,1 \% em termos regionais, obtiveram conceito 3, o conceito modal para a região.
Os conceitos 2 e 4 foram atribuídos a, respectivamente, quatro e 16 cursos ( 8,2 \% e 32,7 \% ).
O conceito 5 foi atribuído a apenas um curso ( 2,0 \% ).
Nessa região tampouco, nenhum dos cursos ficou com conceito 1, enquanto apenas um curso ficou sem conceito ( SC ).
Dos 122 cursos participantes da região Sudeste, como já comentado, 48 ( 39,3 \% ) obtiveram conceito 3, o conceito modal.
O conceito 1 foi atribuído a quatro cursos ( 3,3 \% ) e o conceito 2, a 13 ( 10,7 \% ). 45 cursos ( 36,9 \% ) receberam o conceito 4 e dez outros ( 8,2 \% ) receberam o conceito 5.
Nesta região, apenas dois cursos ficaram sem conceito ( SC ).
A região Sul \textbf{contou} 55 cursos, onde a predominância do conceito 4, como já comentado, foi de 38,2 \%, correspondentes a 21 dos 55 cursos participantes na região Sul.
O conceito 2 foi atribuído a cinco cursos ( 9,1 \% ) e o conceito 3, a 20 cursos ( 36,4 \% ).
O conceito 1 foi recebido por três cursos ( 5,5 \% ) e o conceito 5, por quatro deles ( 7,3 \% ).
Nesta região, apenas dois cursos ficaram sem conceito ( SC ).
Exatamente metade dos 28 cursos participantes na região Centro-Oeste recebeu conceito 3, o conceito modal.
Os demais cursos foram avaliados com conceito 1 ( um curso, 3,6 \% ), conceito 2 ( seis cursos, 21,4 \% ) e conceito 4 ( sete cursos, 25,0 \% ).
Nesta região nenhum curso ficou sem conceito ( SC ) ou tampouco foi alocado no conceito 5. 5.2 CONCEITOS POR CATEGORIA ADMINISTRATIVA E POR GRANDE REGIÃO A Tabela 5.2 apresenta a distribuição dos cursos participantes do ENADE/2012 de Comunicação Social - Jornalismo, por Categoria Administrativa, de acordo com os conceitos por eles alcançados, segundo as Grandes Regiões brasileiras.
Dos 277 cursos participantes, 58 ( 20,9 \% ) eram \textbf{ministrados} em instituições públicas e 219 ( 79,1 \% ), em privadas.
De acordo com as informações da Tabela 5.2, em termos nacionais, entre instituições públicas e privadas : das 15 IES que receberam o conceito 5, dez eram públicas e cinco privadas.
Dos 58 cursos participantes de IES públicas, o conceito 4 foi o valor modal, atribuído a 22 cursos ( 37,9 \% ).
Entre os demais cursos de instituições públicas participantes, cinco obtiveram conceito 1 ( 8,6 \% da categoria ), sete receberam conceito 2 ( 12,1 \% ) e 14 foram alocados ao conceito 3 ( 24,1 \% ).
Como já comentado, os dez \textbf{restantes} receberam o conceito 5.
Nesta categoria, nenhum dos cursos ficou sem conceito ( SC ).
Na rede privada, o conceito modal foi 3, com 105 cursos dos 219 da categoria.
Nesta categoria administrativa, cinco dos cursos ficaram sem conceito ( SC ).
Entre os demais cursos participantes, quatro receberam conceito 1 e 25, conceito 2.
O conceito 4 foi atribuído a 75 cursos e o conceito 5, como já comentado, a cinco cursos.
Na análise por região, observa -se que, na região Norte, as instituições privadas participaram com 15 cursos ( 65,2 \% do total regional, a menor proporção entre todas as regiões ), dos quais nenhum ficou sem conceito ( SC ).
Nenhum curso foi alocado aos conceitos 1 ou 5.
O conceito modal para as instituições privadas na região foi 3, com oito cursos, correspondendo a 53,3 \%.
Os demais cursos receberam os conceitos 2 ( quatro cursos, 26,7 \% ) e 4 ( três cursos, 20,0 \% ).
As instituições públicas participaram com oito cursos ( 34,8 \% do total regional ), dos quais nenhum ficou sem conceito ( SC ).
O conceito modal para esta Categoria Administrativa, na região Norte, foi 4, com cinco cursos correspondendo a 62,5 \%.
Os demais foram alocados aos conceitos 1 ( um curso, 12,5 \% ) e 3 ( dois cursos, 25,0 \% ).
Na região Nordeste, a rede privada concentrou 34 dos 49 cursos participantes, equivalentes a 69,4 \% do total da região.
Nesta Categoria Administrativa na região, predominaram os cursos com conceito 3, 20 cursos correspondendo a 58,8 \% dos cursos oferecidos por IES privadas no Nordeste.
Apenas um curso ficou sem conceito ( SC ).
Os demais foram avaliados com conceitos 2 ( três cursos correspondendo a 8,8 \% ) e 4 ( dez cursos correspondendo a 29,4 \% ).
Nenhum curso recebeu o conceito 1, nem tampouco o conceito 5.
As instituições públicas dessa região participaram com 15 cursos ( 30,6 \% ), dos quais sete obtiveram conceito 3, também o valor modal.
Os oito demais foram avaliados com os conceitos 2 ( um curso ), 4 ( seis cursos ) e 5 ( também um curso ).
Nenhum curso ficou sem conceito ( SC ) nesta combinação de região e categoria.
Na região Sudeste, a proporção de cursos da rede privada, 86,9 \%, foi mais elevada do que nas demais regiões brasileiras, correspondendo a 106 dos 122 cursos participantes.
Nesta combinação de Categoria Administrativa e Grande Região, apenas dois cursos não receberam conceito ( SC ).
O conceito modal foi 3 ( 46 cursos ).
Os demais foram avaliados com conceito 1 ( três cursos ), conceito 2 ( nove cursos ), conceito 4 ( 42 cursos ) e conceito 5 ( quatro cursos ).
Entre os 16 cursos em instituições públicas na região Sudeste, a categoria modal foi a que recebeu conceito 5 ( seis cursos ).
Os demais foram avaliados com os conceitos 1 ( um curso ), 2 ( quatro cursos ), 3 ( dois cursos ) e conceito 4 ( três cursos ).
Na região Sudeste, nesta categoria, nenhum curso ficou sem conceito ( SC ).
As instituições privadas concentraram 42 dos 55 cursos participantes da região Sul, 76,4 \% do total regional.
Desses, 18 ficaram com conceito 3, o conceito modal.
Nesta combinação de Categoria Administrativa e Grande Região, dois cursos ficaram sem receber conceito ( SC ).
Os demais foram avaliados com conceito 1 ( um curso ), conceito 2 ( quatro cursos ), conceito 4 ( 16 cursos ) e conceito 5 ( também um curso ).
As instituições públicas na região Sul participaram com 13 cursos ( 23,6 \% ), dos quais nenhum ficou sem conceito ( SC ).
O conceito modal foi 4, com cinco cursos.
Os demais foram avaliados com conceito 1 ( dois cursos ), conceito 2 ( um curso ), conceito 3 ( outros dois cursos ) e conceito 5 ( três cursos ).
Na região Centro-Oeste, 22 dos 28 cursos participantes eram de instituições privadas ( 78,6 \% em termos regionais ).
Destes, 13 concentraram -se no conceito 3, conceito modal.
Os demais receberam conceito 2 ( cinco cursos ) e conceito 4 ( quatro cursos ).
Nenhum curso ficou sem conceito ( SC ) nesta categoria, nem tampouco foi alocado aos conceitos 1 ou 5.
Dos seis cursos de instituições públicas, três foram avaliados no conceito 4, conceito modal, e os outros três cursos foram alocados igualmente aos conceitos 1, 2 e 3.
Nesta região, nenhum curso de IES \textbf{pública} ficou sem conceito ( SC ). 5.3 CONCEITOS POR ORGANIZAÇÃO ACADÊMICA E POR GRANDE REGIÃO Na Tabela 5.3 encontra -se a distribuição dos conceitos atribuídos aos cursos participantes do ENADE/2012 na Área de Comunicação Social - Jornalismo, por Organização Acadêmica, segundo as Grandes Regiões brasileiras.
Dos 277 cursos participantes, 130 eram oferecidos em Universidades, 49 em Centros Universitários e os demais 98 em Faculdades.
Esta distribuição corresponde a, respectivamente, 46,9 \%, 17,7 \% e 35,4 \% dos cursos.
De acordo com os dados apresentados, de todos os 15 cursos avaliados com conceito 5, dez eram vinculados a Universidades.
Apenas um curso em Universidades ficou sem conceito ( SC ).
Esse tipo de Organização Acadêmica teve o conceito 4 como modal, com 54 cursos ( 41,5 \% ).
Os demais cursos de Universidades avaliados receberam os conceitos 1 ( seis cursos ), 2 ( dez cursos ), 3 ( 49 cursos ) e conceito 5 ( dez cursos, como já mencionado ).
Entre os cursos em Centros Universitários, o conceito modal foi 3, com 19 cursos ( 38,8 \% ).
Neste tipo de organização acadêmica nenhum curso ficou sem conceito ( SC ).
Os outros cursos neste tipo de Organização Acadêmica receberam os conceitos 1 ( um curso ), 2 ( dez cursos ), conceito 4 ( 17 cursos ) e conceito 5 ( dois cursos ).
Nas Faculdades, quatro dos 98 cursos ficaram sem conceito ( SC ) e 51 ( 52,0 \% ) receberam o mesmo conceito modal dos Centros Universitários, conceito 3.
Dos demais cursos neste tipo de Organização Acadêmica, dois receberam conceito 1, 12 conceito 2, 26 cursos conceito 4 e três conceito 5.
Considerando -se separadamente as regiões brasileiras, verifica -se que, na região Norte, as Universidades concentraram nove dos 23 cursos participantes.
Dos cursos em Universidades, nenhum ficou sem conceito ( SC ) ou recebeu conceitos 2 ou 5.
Cinco cursos receberam o conceito modal, 4.
Os demais quatro ficaram distribuídos, um curso no conceito 1 e três cursos no conceito 3.
Os Centros Universitários da região Norte foram representados por quatro cursos, nenhum sem conceito ( SC ).
Três receberam conceito 3, o conceito modal, e um recebeu o conceito 4.
As Faculdades participaram com dez cursos na região Norte.
Nenhum ficou sem conceito ( SC ) ou recebeu os conceitos 1 ou 5.
Os demais receberam conceito 2 ( quatro cursos ), 3 ( outros quatro cursos ) e conceito 4 ( dois cursos ).
Na região Nordeste, as Universidades participaram com 19 dos 49 cursos na Área de Comunicação Social - Jornalismo da região.
Nenhum dos cursos em Universidades no Nordeste ficou sem conceito ( SC ) ou recebeu conceito 1.
O conceito modal foi 3, com nove cursos.
Os demais receberam conceito 2 ( um curso ), conceito 4 ( oito cursos ) e conceito 5 ( também um curso ).
Os Centros Universitários \textbf{contaram} com quatro cursos participantes na região Nordeste, dois deles com conceito 3 e os outros dois com conceitos 2 e 4, um em cada.
As Faculdades foram representadas por 26 cursos na região Nordeste.
O conceito modal foi 3, com 16 cursos.
Dos cursos \textbf{restantes}, dois receberam conceito 2, sete receberam conceito 4, enquanto apenas um curso ficou sem conceito ( SC ).
Na região Sudeste, as Universidades concentraram 59 dos 122 cursos da região.
Entre os cursos em Universidades na região, os conceitos modais foram 3 e 4, com 22 cursos cada, enquanto apenas um curso ficou sem conceito ( SC ).
Os demais cursos receberam os conceitos 1 ( dois cursos ), 2 ( seis cursos ) e 5 ( outros seis cursos ).
Os Centros Universitários participaram com 30 cursos na região Sudeste, dos quais 12 obtiveram conceito modal, 4, e nenhum ficou sem conceito ( SC ).
Os demais receberam os conceitos 1 ( um curso ), 2 ( cinco cursos ), conceito 3 ( 11 cursos ) e conceito 5 ( também um curso ).
As Faculdades foram representadas por 33 cursos na região Sudeste, dos quais 15 obtiveram conceito modal 3.
Os demais receberam os conceitos 1 ( um curso ), 2 ( dois cursos ), conceito 4 ( 11 cursos ), conceito 5 ( três cursos ), enquanto apenas um curso ficou sem conceito ( SC ).
Dos 55 cursos da região Sul, 31 eram de Universidades, para os quais o conceito modal foi 4, com 14 cursos.
Nesse tipo de organização, nenhum dos cursos ficou sem conceito ( SC ) e os demais receberam os conceitos 1 ( dois cursos ), 2 ( um curso ), 3 ( 11 cursos ) e 5 ( três cursos ).
Os Centros Universitários da região Sul tiveram três dos sete cursos participantes no conceito modal, 2.
Os outros quatro cursos receberam conceitos 3 ( um curso ), 4 ( dois cursos ) e 5 ( também um curso ).
Foram 17 os cursos vinculados a Faculdades na região Sul, e apenas dois destes ficaram sem conceito ( SC ).
Oito desses cursos receberam conceito 3, o conceito modal.
Os demais receberam os conceitos 1 ( um curso ), 2 ( também um curso ) e 4 ( cinco cursos ).
Na região Centro-Oeste, 12 dos 28 cursos eram de Universidades.
Nesse tipo de organização, nenhum curso ficou sem conceito ( SC ) ou recebeu o conceito 5.
O conceito modal foi 4, com cinco cursos.
Os outros cursos obtiveram os conceitos 1 ( um curso ), 2 ( dois cursos ) e 3 ( quatro cursos ).
Os Centros Universitários da região Centro-Oeste \textbf{contaram} com quatro cursos, dos quais dois receberam o conceito 3 e os outros dois os conceitos 2 e 4 ( um cada ).
Nenhum curso ficou sem conceito ( SC ).
Dos 12 cursos em Faculdades na região Centro-Oeste, nenhum ficou sem conceito ( SC ), três receberam conceito 2, oito conceito 3 e o outro curso conceito 4.

% matched lemmas: contar, equivaler, ministrar, panorama, pública, restante, tabular
\end{textsample}
